% Transcription — fonds Alexandre Grothendieck, shelfmark 67, batch 4 (pages 61-80).
% Established from the facsimile archives/batches/67/batch-04.pdf (a mirror of
% https://grothendieck.umontpellier.fr/67.pdf), per the skill
% transcribe-grothendieck. The batch file carries no cover sheet: PDF page N
% is page 60+N of the folder. The Montpellier footer printed under the leaf
% confirms it (« 65 » under PDF page 5, « 70 » under PDF page 10).
%
% Pass: Opus 5.5 (claude-opus-5-5), 2026-09-22 — first pass, unchecked against the pages by a human.
%
% All twenty pages are mathematics and all are transcribed; nothing skipped.
% \page{} numbers the position in the folder, as folder 1 does: the
% archivists' pencil numbers bottom-left run three ahead here (page 61 is
% pencilled 64, page 80 is pencilled 83).
%
% What the batch is. One continuous run of ink notes, rectos numbered by him
% 5 to 14 (pages 61, 63, …, 79; the versos 62, 64, …, 80 carry the
% continuation and no number of his). It begins mid-sentence: his pp. 1-4
% fall in batch 3.
%   61-69  Discs in a conformal surface X' around marked points, their
%          « ombres » (shadows) in the tangent spaces T_{X',a}; the
%          admissibility condition (the canonical constant-curvature metric
%          on X* = X minus the open discs, boundary geodesic, the two toric
%          structures on each ∂D_i agreeing); the Teichmüller space with
%          boundary TB_{g,ν} → T_{g,ν}; a « Conjecture » (page 67) that
%          admissible disc systems ↔ shadow systems is bijective, so that
%          TB_{g,I} → T_{g,I} is a principal (R*+)^I-bundle; unease at the
%          implied canonical Euclidean metrics on T_{X,a_i} when g = 0.
%   69-73  Dated « 12.1. » in the left margin: « 5. Retour sur le principe
%          de chirurgie conforme » — Lemme 4 A) and B): uniformising a
%          Jordan domain slit along an arc I by the closed disc, the
%          rectifiable gluing homeomorphism φ : β → β̄, and a « Scholie »
%          (bijection between rectifiable homeos of [0,1], rectifiable slits
%          I in D, maps f : D → D, power series).
%   73-80  Dated « 13.1. », section ⑥: from a boundary Teichmüller object
%          (X, C) to a stable curve X''' with normal propagation structure
%          and admissible discs; the moduli-level maps MDB_{g,ν} →
%          M_stND_{g,ν} over mdt_{g,ν}, and MD_{g,ν} → M_stN_{g,ν} over
%          md_{g,ν}; « Question » whether they are isomorphisms, the key
%          case (g,ν) = (0,4) (two Thurston pants glued), and the effect on
%          fundamental groups (Teichmüller groups). The batch ends
%          mid-sentence (« de vérifier si »).
% The two margin dates « 12.1. » and « 13.1. » carry no year; they are
% transcribed as written and not interpreted.
%
% Notation fixed for the batch: D, D_i discs, a, a_i their centres, o(D)
% « ombre de D », T_{X',a} tangent space, X* = X \ ∪ D̊_i, \mathbb{D} the
% closed unit disc and \mathring{\mathbb{D}} its interior, ζ = exp 2iπ/3,
% α, β, β̄ the three arcs; MDB, M_{st}ND, mdt, MD, M_{st}N, md set in
% \mathrm with his subscripts g,ν; X' the cut surface, X'' the capped
% compact surface, X''' the stable curve; T^* the C*-torsor of non-zero
% vectors (his star), ∧_{C*} the contracted product.
%
% The dating is the inventory's own for the folder, copied verbatim. Its
% group, [66-89] « Autour de l'enseignement - Thèmes liés (cours et
% directions de recherches) », is dated 1964-1984.
\documentclass[11pt,a4paper]{article}
% Le préambule commun à toutes les transcriptions du fonds Grothendieck.
%
% Il tient en une idée : **l'incertitude se compose, elle ne se résout pas.**
% Une transcription qui lisse un mot douteux a détruit la seule chose pour
% laquelle on la faisait — un lecteur doit pouvoir savoir, à chaque endroit,
% ce qui était sur le feuillet et ce qui a été ajouté. Les six macros qui
% suivent sont donc l'appareil critique, et non de la décoration.
%
% Les mêmes commandes sont reconnues par `scripts/render.mjs`, qui produit la
% vue de lecture du site. Une macro ajoutée ici doit l'être aussi là-bas,
% faute de quoi le rendu échouera bruyamment — ce qui est voulu : un rendu
% silencieusement faux est pire qu'un rendu absent.

\ProvidesPackage{grothendieck}[2026/08/08 Transcriptions du fonds Grothendieck]

\usepackage{amsmath,amssymb,amsthm}
\usepackage{mathrsfs}
\usepackage{stmaryrd}
% Les diagrammes commutatifs sont partout dans ce fonds : c'est la langue
% même de la géométrie algébrique de Grothendieck, et une transcription qui
% les rendrait en prose ne transcrirait rien.
\usepackage{tikz-cd}
% Français par défaut — c'est la langue des feuillets — mais l'anglais est
% chargé aussi : la traduction et le résumé sont des documents anglais, et la
% typographie française (espaces avant les deux-points) y serait une faute.
% Ces documents basculent par \selectlanguage{english} en tête de corps.
\usepackage[english,french]{babel}
\usepackage{fontspec}
\usepackage{geometry}
\geometry{a4paper,margin=25mm,marginparwidth=28mm}
\usepackage{marginnote}
\usepackage{xcolor}
% ulem, not soul. `soul`'s \st and \ul are built on a character-by-character
% scan that XeLaTeX's font machinery breaks: the document fails with an
% undefined control sequence rather than degrading. ulem does the same two jobs
% and is Unicode-engine safe. `normalem` stops it hijacking \emph, which the
% transcriptions use for Grothendieck's own emphasis.
\usepackage[normalem]{ulem}
\usepackage{hyperref}
\hypersetup{colorlinks=true,linkcolor=black,urlcolor=black,pdfborder={0 0 0}}

% --- Métadonnées du lot -----------------------------------------------------
% Reprises telles quelles par le script de rendu, qui les affiche en tête de
% la vue de lecture. Elles fixent ce que le fichier prétend couvrir : sans
% elles, une transcription flotte sans point d'attache dans un fonds de seize
% mille feuillets.

\newcommand{\folder}[1]{\gdef\grothFolder{#1}}
\newcommand{\batch}[1]{\gdef\grothBatch{#1}}
\newcommand{\leaves}[2]{\gdef\grothFirst{#1}\gdef\grothLast{#2}}
\newcommand{\foldertitle}[1]{\gdef\grothFoldertitle{#1}}
\gdef\grothFolder{?}\gdef\grothBatch{?}\gdef\grothFirst{?}\gdef\grothLast{?}\gdef\grothFoldertitle{}

% --- Le résumé --------------------------------------------------------------
%
% Il ouvre la lecture modernisée, sous le titre « Résumé », et il est composé
% en petit. Ce n'est pas un ornement : c'est le seul passage du document qui
% ne soit pas des mathématiques, et un lecteur doit pouvoir le reconnaître
% avant de l'avoir lu — puis le sauter s'il connaît déjà le sujet, ou s'y
% arrêter s'il ne le connaît pas. Le filet qui le suit marque la reprise du
% corps.

\newenvironment{resume}
  {\small\setlength{\parskip}{0.45em}}
  {\par\medskip\hrule\medskip}

% --- Mots-clés ---------------------------------------------------------------
%
% En anglais, en clôture du résumé de la lecture modernisée : le vocabulaire
% moderne sous lequel chercher ce que les feuillets construisent. Le manifeste
% du site (`npm run manifest`) les extrait pour étiqueter la cote entière —
% cette ligne est la source unique des tags, il n'y a pas de fichier à côté.
\newcommand{\keywords}[1]{\par\smallskip\noindent
  {\footnotesize\textsc{Keywords} --- \textit{#1}}\par}

% --- L'appareil critique ----------------------------------------------------

% Le numéro de feuillet, en marge. C'est l'ancre qui permet de régler un
% désaccord sur une formule en regardant *un* feuillet plutôt qu'une chemise
% de six cents. Le site s'en sert aussi pour tourner les pages du fac-similé
% quand on fait défiler la transcription.
\newcommand{\leaf}[1]{%
  \marginnote{\footnotesize\color{gray!70}\textsf{#1}}%
  \label{leaf:#1}\ignorespaces}

% Illisible : on ne devine pas. Un mot inventé qui se lit comme les autres est
% le pire résultat possible.
\newcommand{\ill}{\textcolor{red!60!black}{[\,\ldots\,]}}

% Lecture douteuse : le mot est donné, et signalé comme incertain.
\newcommand{\uncertain}[1]{\uline{#1}}

% Ajout de l'éditeur — une lettre manquante, un mot sous-entendu.
\newcommand{\add}[1]{\textcolor{blue!50!black}{[#1]}}

% Biffé par Grothendieck lui-même. Ce qu'il a rayé fait partie du document :
% une rature dit souvent où la pensée a changé de direction.
\newcommand{\struck}[1]{\sout{#1}}

% Note du transcripteur, sur l'état matériel du feuillet.
\newcommand{\note}[1]{\par\smallskip{\small\color{gray!60!black}\textit{[#1]}}\par\smallskip}

% Note marginale de Grothendieck, distincte du corps du texte.
\newcommand{\marginal}[1]{\par\smallskip\noindent\hspace{1em}%
  {\small\color{gray!60!black}#1}\par\smallskip}

% --- Titre ------------------------------------------------------------------

\AtBeginDocument{%
  \begin{center}
    {\large\bfseries Fonds Alexandre Grothendieck --- cote n\textsuperscript{o}~\grothFolder}\\[2pt]
    {\itshape\grothFoldertitle}\\[4pt]
    {\small feuillets \grothFirst--\grothLast\ (lot \grothBatch)}\\[2pt]
    {\footnotesize\color{gray!60!black}Transcription établie d'après le fac-similé numérisé
     par l'Université de Montpellier.\\
     Les lectures incertaines sont soulignées, les illisibles notées [\,\ldots\,],
     les ajouts entre crochets.}
  \end{center}
  \vspace{1em}}

\endinput


\folder{67}
\batch{4}
\pages{61}{80}
\dating{[à partir de 1977]-1984}
\watermark{Édition de démonstration}
\foldertitle{Chirurgie des surfaces conformes : notes manuscrites (s.d., 1983-1984), lettres (1984)}

\begin{document}

\page{61}
\note{pagination de l'auteur : 5. La page commence au milieu d'une phrase
commencée au lot précédent.}

les propriétés de régularité de $C$ \uncertain{au fait que} \uncertain{cercle}
« \uncertain{intérieur} » dans la surface \uncertain{conforme}
\struck{\ill{}} $X'$ déduite par bordage des trous. D'autre part, le choix
d'un \uncertain{point} de $D \setminus C$ équivaut à celui d'une structure
torique sur $C$ compatible avec la structure de droite projective réelle
(liée aux restrictions du groupe structural
$\mathrm{GL}(2,\mathbb{R})/\mathbb{R}^{*}$ ($= \mathrm{GP}(1,\mathbb{R})$)
au sous-groupe compact maximal $\mathrm{O}(2)/\pm 1$) -- et est équivalent
aussi à la donnée d'une structure \uncertain{torique} de $D$ sur le bord
$\partial D = C$. Notons que l'on aura alors une inclusion canonique
\[
  D \subset T_{X',a}
\]
où $a$ est le centre choisi de $D$. On va expliciter \struck{ainsi} ces
dernières relations ainsi.
\note{le signe entre $D$ et $C$ est un petit trait, lu comme une
différence ; « $= \mathrm{GP}(1,\mathbb{R})$ » est écrit au-dessus de
$\mathrm{GL}(2,\mathbb{R})/\mathbb{R}^{*}$ avec une accolade ; « le bord »
est en interligne.}

\emph{Prop.} Équivalence entre \struck{la cat.\ des surfaces conf.}

\page{62}

a) Cat.\ des surfaces conformes à bord, \uncertain{compactes}, avec comp.\
connexe choisie $C$ du bord, et une structure torique disons rectifiable
(?) pour la structure analytique réelle, et

b) Cat.\ des surfaces conformes \struck{bord} $X'$, avec \struck{\ill{}}
disque pointé $(D, a)$, $D \subset \operatorname{Int}(X')$ \ill{}, tel que
$\partial D$ soit rectifiable dans $X'$ (i.e.\ l'injection
$\partial D \to X'$ rectifiable, pour la structure analytique naturelle de
$\partial D$ comme bord de $D$).

De plus, pour $(X', a)$ fixé, chaque disque $D$ autour de $a$ définit un
« disque \struck{canonique} euclidien » dans $T_{X',a}$ (\ill{} la structure
euclidienne (de $T_{X',a}$) \uncertain{conforme}). Comme les disques forment
un torseur sous $\mathbb{R}^{*+}$ (par homothéties), \uncertain{i.e.} un
espace de dim 1, il est clair que l'application

\page{63}
\note{pagination de l'auteur : 6.}

\[
  D \longmapsto o(D) \qquad (\text{« ombre de } D \text{ »})
\]
qui s'ensuit (des disques conformes \uncertain{centrés} en $a$)
$D \mapsto$ ombre \struck{de} dans $T_{X',a}$, n'a rien d'injectif. Les
\emph{disques « concentriques »} à $D$ donnent les $\lambda\, o(D)$
($0 < \lambda \leq 1$), plus précisément, $D \mapsto o(D)$ est compatible
aux homothéties de rapport \struck{$\lambda \leq$} $\lambda \in {]0,1]}$.

Je m'intéresse \uncertain{ici} au cas où $X'$ est compacte, et munie d'une
famille finie de « pts marqués », soit $(a_i)_{i \in I}$, \struck{$X$}, et
où on prend une famille \uncertain{mutuellement disjointe} $D_i$
($i \in I$) autour des $a_i$ de disques \uncertain{donnant} des ombres
\[
  o(D_i) \subset T_{X',a_i} .
\]
On demande des conditions restrictives \uncertain{plausibles} sur le
système des $D_i$, impliquant (qui sait ?) qu'il est déterminé par le
système des ombres.

\page{64}

[On laisse tomber dorénavant le $'$ des $X'$ !!]

\struck{Mais sur $X^{*} = X \setminus \bigcup D_i$ (si chaque composante)}
\struck{\ill{}} \uncertain{OPS} $X$ connexe, $I$ non vide -- \struck{dans
$X$ \ill{}} il y a sur $X^{*} = X \setminus \bigcup_i \mathring{D}_i$
(surface conforme : bord contenu dans le disque de $X$ \ill{}
$\bigcup \partial D_i$) une \struck{\ill{}} métrique
\struck{hyperbolique} \uncertain{canonique} à courbure constante (qu'on peut d'ailleurs
normaliser par la condition que la \uncertain{courbure} \uncertain{soit}
$\pm 1$) \struck{sont des} \uncertain{sauf cts} canonique dans le cas
hyperbolique (i.e.\ \uncertain{hors du cas} $g = 0$,
$\nu \in \{1, 2\}$, $\nu = \operatorname{card} I$) -- dans le cas
hyperbolique les bords sont géodésiques, dans tous les cas le bord est muni
d'une structure torique canonique provenant de la structure riemannienne.
Il est \uncertain{tentant} de restreindre les \add{syst.\ de} disques
envisagés, en exigeant que les deux structures toriques sur
\struck{\ill{}} chaque $\partial D_i$ (provenant soit de $D_i$, soit de
$X^{*}$)
\marginal{il y a aussi le cas \struck{deux} où \ill{} du plan projectif
\struck{\ill{}} réel ($\simeq S^2/\pm 1$) \struck{d'un disque} « un trou »}
\note{« syst.\ de » est en interligne. La note marginale est écrite en
biais dans la marge gauche, au niveau de la condition sur $g$ et $\nu$ ;
elle se lit mal. « OPS » ouvre la ligne, souligné, au-dessus d'un mot
biffé ; la lecture « $\pm 1$ » de la normalisation est incertaine.}

\page{65}
\note{pagination de l'auteur : 7.}

coïncident. Cela signifie donc que $X$ se déduit de $X^{*}$ par « bordage
des trous par des disques canoniques » (utilisant les structures toriques
sur les comp.\ conn.\ de $\partial X^{*}$). \struck{La question}
\uncertain{Disons alors} \struck{\ill{}} que le système de disques
$(D_i)_{i \in I}$ dans $X$ autour des $a_i$ est « admissible ».

\emph{Questions} Comment cette condition d'admissibilité se reflète-t-elle
en termes de la géométrie intrinsèque de $X$ ? Par exemple, \struck{quand}
quand $X$ est muni d'une métrique \ill{} canonique (cas où
$X \not\simeq \mathbb{P}^1_{\mathbb{C}}$), quelle est la relation de cette
condition avec celle que les $\partial D_i$ soient des « disques
géodésiques » (définis par la condition que
$\mathrm{dist.\,géod}(x, a_i) \leq r_i$) ?

Fixons une $X_0$ surface top.\ \uncertain{orientée} \struck{à bord} de
référence ($X_0$, de type $g,\nu$,
\marginal{\ill{} type $(g,\nu)$ \ill{} $(0,1)$, $(0,2)$ \ill{}}
\note{la note marginale, au bas de la marge gauche, est très raturée ;
seuls les couples $(0,1)$, $(0,2)$ et « type $(g,\nu)$ » se lisent.}

\page{66}

et considérons « l'espace de Teichmüller à bord » correspondant
\struck{$T_{g,\nu}$} $\mathrm{TB}_{g,\nu}$, qui est \struck{isom.}\
\uncertain{homéo.}\ à $\mathbb{R}^{6g-6+3\nu}$ ou mieux, à
$\mathbb{C}^{3g-3+\nu} \times \mathbb{R}^{\nu}$. Considérons l'application
\[
  \mathrm{TB}_{g,\nu} \longrightarrow T_{g,\nu}
\]
\note{sous la source est écrit $\simeq \mathbb{C}^{3g-3+\nu} \times
\mathbb{R}^{\nu}$, sous le but $\mathbb{C}^{3g-3+\nu}$, suivi d'un
$\times \mathbb{R}^{\nu}$ biffé.}
de bordage des trous, je présume qu'elle est fibrante de fibre
\uncertain{\ill{}}, plus précisément, que \struck{deux pts de}
\uncertain{c'est la projection} canonique. \struck{et} En effet, cela doit
provenir du fait que pour une sphère \uncertain{holomorphe} à bord de type
$(0,3)$, \struck{\ill{}} quand la \ill{} \uncertain{combinatoire}
(\uncertain{triangle} combinatoire) est fixée, et qu'on fait varier
seulement \struck{les} un, deux ou trois des paramètres « longueur »
correspondant aux trois trous, la sphère « à bord et paramétrisée »
(\uncertain{mono}, bi ou triparamétrée suivant les cas)
\struck{\ill{}} \uncertain{congruente}

\page{67}
\note{pagination de l'auteur : 8.}

par bordage des trous correspondants reste la même (à isom.\ can.\ près).
Il faudrait, pour le voir clairement, voir la relation entre la sphère
\add{à bord} $S(I)$ triparamétrée, associée à la sphère à bords
$\Sigma^{*}$, et ses paramètres $l_i$ ($i \in I$) comme diamètres
(\uncertain{riemanniens}) des $S(I)$ en enlevant des disques qui
\uncertain{dépendent si} ?) ($D(l_i)$ \struck{qui} dont chacun ne dépend
que du $l_i$ correspondant, pas des autres $l_j$.
\note{« à bord » est en interligne ; un grand crochet en marge gauche
embrasse ce passage.}

Ce point étant supposé acquis, ou quelque variante, il semble donc bien que
(au moins dans le cas hyperbolique pointé) pour une surface conforme
pointée donnée, \add{de type $(g,\nu)$,} la famille $(D_i)$ (des disques
autour des $a_i$) forme un espace homéomorphe à $\mathbb{R}^{\nu}$
($\nu = \operatorname{card} I$), i.e.\ à celui des systèmes d'ombres
possibles ! Il est donc très \uncertain{tentant} de faire la

\emph{Conjecture} Pour une surface conforme pointée, (au moins dans le cas
hyperbolique), l'application \struck{que}
\[
  (D_i) \longmapsto (o(D_i))
\]

\page{68}

des syst.\ admissibles de disques \uncertain{conformes} autour des $a_i$,
vers \uncertain{l'ens.\ des} syst.\ de disques \struck{conformes}
\uncertain{euclidiens} dans les $T_{X,a_i}$, est bijective.
\struck{Au niveau des esp…}
Au niveau des espaces modulaires du type Teichmüller, cela donnerait donc
\struck{une} \struck{la} structure de fibré principal de groupe structural
$(\mathbb{R}^{*+})^{I}$ :
\[
  \mathrm{TB}_{g,I} \longrightarrow T_{g,I} \qquad (\operatorname{card} I = \nu),
\]
\note{au-dessus, une première écriture biffée :
$\mathrm{TB}_{g,I} \xrightarrow{\ \sim\ } T_{g,I} \times$ \ill{}.}
Si la \ill{} relation \uncertain{qui} précédait \struck{\ill{}} était vrai,
cela signifierait qu'il y aurait une section canonique de ce fibré, -- ce
qui signifie aussi que pour un pt de $\mathrm{TB}_{g,I}$, i.e.\ une surface
à bord ($X_0$-rigidifiée) de type $(g,\nu)$, il y aurait un syst.\
d'ombres canonique dans les $T_{X,a_i}$, i.e.\ des

\page{69}
\note{pagination de l'auteur : 9.}

métriques euclidiennes \emph{canoniques} sur les $T_{X,a_i}$, décrivant leur
structure conforme. Dans le cas $g \geq 2$ et même pour $g = 1$, cela ne
semble pas inquiétant, puisqu'on a une métrique canonique sur $X$. Par
contre dans le cas $g = 0$, cela semble étrange, puisque dans ce cas on n'a
pas de métrique canonique sur $X$.

\section*{5. Retour sur le principe de chirurgie conforme}
\marginal{12.1.}
\note{le titre est de sa main, souligné ; « 12.1. », souligné, est écrit
dans la marge gauche à sa hauteur, sans année.}

La formulation des lemmes \add{locaux} 2 et 3, qui sont l'ingrédient
technique essentiel pour le Th.\ 4 (et lui sont équivalents), est un peu
ambiguë. Il est plus élégant de se réduire à des lemmes globaux du type
uniformisation par le disque $\mathbb{D}$. Cela donne le

\page{70}

\marginal{Lemme 4}
\marginal{NB Il y a une version plus faible de A), \ill{}
$\dot U = C$ (sans $I$), \ill{} formellement \ill{} A)}
\note{En haut à gauche, une figure : une courbe fermée $C$, un arc $I$
joignant un point $b$ de $C$ à un point $a$ intérieur, et deux branches
marquées $\zeta$, $\bar\zeta$ partant de $b$ vers l'extérieur ; une flèche
indique l'orientation. La note « NB … » est écrite en biais au-dessus de la
figure.}

A) Soit $U$ une partie ouverte bornée de $\mathbb{C}$, dont la frontière
est \supplied{réunion d'}une courbe de Jordan \struck{rectifiable. Alors il
existe une fonction, continue} $C$, et d'un intervalle
\struck{homéomorphe} topologique $I$ tel que \struck{$f : \mathbb{D}$}
$I \cap C = \{b\}$, \struck{\ill{}} (alors $I \setminus \{b\}$ est contenu
dans la comp.\ connexe bornée $V$ de $\mathbb{C} \setminus C$, et $b$ est
une extrémité de $I$, ($U = V \setminus I'$, où $I' = I \setminus \{b\}$)).
\struck{Soit} Supposons de plus que $C$ et $I$ soient rectifiables. Soit
$a$ l'autre extrémité de $I$. Alors il existe une fonction
\struck{continue} $f : \mathbb{D} \to \mathbb{C}$ telle que

1) $f$ est continue, et induit une \struck{ho} bijection (voire un homéo !)
$\mathring{\mathbb{D}} \xrightarrow{\ \sim\ } U$.

2) $f(\zeta) = f(\bar\zeta) = b$, et $f$ induit une bijection (voire un
homéo) $\alpha \setminus \{\zeta, \bar\zeta\} \xrightarrow{\ \sim\ }
C \setminus \{b\}$.
(NB $\zeta = \exp 2i\pi/3$, et $\alpha$, $\beta$, $\bar\beta$ sont les
trois arcs sur $\partial\mathbb{D}$ déterminés par $1, \zeta, \bar\zeta$).
\note{en marge, un petit cercle montre les trois arcs : $\alpha$ du côté
de $\zeta$ et $\bar\zeta$, $\beta$ entre $1$ et $\zeta$, $\bar\beta$ entre
$1$ et $\bar\zeta$.}

3) $f|\beta$ et $f|\bar\beta$ induisent des bijections de $\beta$ et
$\bar\beta$ sur $I$).

NB a) b) c) peuvent s'exprimer en disant que $f$ induit un \struck{\ill{}}
\ill{} de $\mathbb{D}$ sur \ill{} déduit de $\overline{U} = \overline{V}$
suivant $I$ (on suppose $I \cup C = \dot U$).
\note{les conditions sont numérotées 1) à 4) et renvoyées ici comme a),
b), c) : ainsi sur la page.}

4) $f|\mathring{\mathbb{D}}$ est holomorphe.

\page{71}
\note{pagination de l'auteur : 10.}

\emph{NB 1} Cette fonction est unique, \uncertain{dû au} fait que
\struck{le} groupe des automorphismes \struck{conformes} holomorphes de
$\mathbb{D}$ agit de façon simplement transitive sur l'ens.\ des triples
\ill{} de pts distincts $(r, s, t)$ \add{de} $S^1 = \partial\mathbb{D}$ qui
se suivent pour l'orientation canonique de \struck{\ill{}} $S^1$.
\struck{\ill{}}

\emph{NB 2} La condition 3) implique qu'on a une \struck{\ill{}} bijection
canonique
\[
  \varphi : \beta = \widehat{(1,\zeta)} \xrightarrow{\ \sim\ }
  \bar\beta = \widehat{(1,\bar\zeta)}
\]
telle que $\varphi(1) = 1$, $\varphi(\zeta) = \bar\zeta$).

B) Cette application $\varphi$ est rectifiable. Inversement,
\struck{tout} \struck{application} \uncertain{homéo.}\ rectifiable
$\varphi : \beta \xrightarrow{\ \sim\ } \bar\beta$ provient d'une situation,
comme dans A. De façon plus précise, \struck{il existe, \ill{}} (quitte à
remplacer $\overline{V}$ par $\mathbb{D}$ par la th.\ d'uniformisation
\uncertain{A) simplifiée} sans segment intérieur $I$, et normaliser la
situation en \struck{\ill{}} amenant $a$ en $0$, $b$ en $1$)
\struck{il existe une} Il existe une application continue
\[
  f = f_{\varphi} : \mathbb{D} \longrightarrow \mathbb{D}
\]
telle que

1) $f|\mathring{\mathbb{D}}$ soit \struck{holo} injection holomorphe

\page{72}

\note{En haut à gauche, une figure : le disque, avec un arc ondulé
joignant le centre $0$ au point $1$ du bord.}

(soit $U = f(\mathring{\mathbb{D}})$)

2) $f(0) = 0$, $f(\zeta) = f(\bar\zeta) = 1$

\struck{3) $f|$ \ill{}} $f|(\alpha \setminus \{\zeta, \bar\zeta\})$
\struck{est un} induit un homéomorphisme avec
$\partial\mathbb{D} \setminus \{1\}$

3) \struck{$f|[\zeta, \bar\zeta]$} $f|\beta$ et $f|\bar\beta$ sont
\uncertain{bijectives} \ill{} \struck{est} $=$ image \uncertain{$I$},
\struck{qui contient \ill{} et $I \setminus \{1\}$} d'après
2) et $I \setminus \{1\}$ est contenu dans $\mathring{\mathbb{D}}$.
\note{ce passage est corrigé en place et en interligne ; l'ordre des mots
rendu ici est incertain.}

4) Soient $f_1 : \beta \xrightarrow{\ \sim\ } I$,
$f_2 : \bar\beta \xrightarrow{\ \sim\ } I$ induits par $f$, alors
$\varphi = f_2^{-1} f_1$.

\emph{NB} Si on veut, $f$ est unique. On peut se servir du théorème
d'uniformisation avec bord dans le \struck{Scholie}

\emph{Scholie} Il y a corr.\ 1-1 entre les données suivantes :

a) Les homéos rectifiables \struck{\ill{}}
\[
  \psi : [0,1] \xrightarrow{\ \sim\ } [0,1]
\]
avec $\psi(0) = 0$ (donc $\psi(1) = 1$)

b) Les segments topologiques \emph{rectifiables} $I$ dans $\mathbb{D}$,
tels que \struck{$I \cap \partial\mathbb{D}$ soit réduit} $I$
\uncertain{admette} $0, 1$ comme extrémités et \struck{\ill{} à une
extrémité de $I$} $I \cap \partial\mathbb{D} = \{1\}$ (i.e.\
$I \setminus \{1\} \subset \mathring{\mathbb{D}}$)

c) Les \struck{\ill{}} applications continues
$f : \mathbb{D} \to \mathbb{D}$ qui satisfont les propriétés 1) 2) 3) de B)
ci-dessus

\page{73}
\note{pagination de l'auteur : 11.}

d) Les séries convergentes
\[
  f = \sum_{i \geq 0} a_i z^i \qquad (a_i \in \mathbb{C})
\]
de rayon de convergence $\geq 1$ (donc définissant
$\mathring{\mathbb{D}} \to \mathbb{C}$) \add{avec $|f(z)| \leq 1$ si
$z \in \mathbb{D}$}, telles que la $f$ se prolonge par continuité en une
fonction $\mathbb{D} \to \mathbb{D}$, celle-ci ayant les propriétés 1) 2)
3) de B) ci-dessus.

\section*{6}
\marginal{13.1.}
\note{le « 6 » est cerclé ; « 13.1. », souligné, est dans la marge gauche,
sans année, au lendemain du « 12.1. » de la page 69.}

Soit $(X, C)$ un objet du groupoïde de Teichmüller analytique complexe à
bord, avec la condition \uncertain{différentiable}. Pour fixer les idées,
supposons $X$ connexe, de type $g,\nu$. Considérons la décomposition $X'$
de $X$ suivant $C$, bouchons ses trous (canoniquement en utilisant les
structures toriques de $\partial X'$)$^{*}$, d'où $X''$ surface holomorphe
compacte \uncertain{sans bord}, à disques pointés marqués, avec des
\struck{iso} \uncertain{antiiso} \uncertain{unitaires} \struck{de ces} de
ces disques (qui se groupent \struck{deux à deux} par paires, à
l'exception des $\nu$ disques correspondant à $\partial X$).
\marginal{$^{*}$ via la métrique riemannienne can.\ de $X'$}
On peut considérer $X''$ comme la normalisée d'une courbe
\uncertain{$\mathbb{C}$-alg.}\ $X'''$ \struck{\ill{}} « stable » (au sens
de Mumford-Deligne), en

\page{74}

identifiant les pts des centres des disques qui se correspondent, et en
pointant par les centres des autres disques. De plus, \uncertain{on}
obtient, \struck{\ill{}} pour chaque point singulier $x$ de $X'''$,
correspondant à une paire $\{x_1, x_2\}$ de points de $X''$ (centres des
deux disques associés), grâce à l'antiisomorphisme des \uncertain{mêmes}
disques pointés, un élément \struck{\ill{}} (i.e.\ \ill{}) de
\[
  T_{X'',x_1} \otimes_{\mathbb{C}} T_{X'',x_2}
\]
i.e.\ un élément de
\[
  T^{*}_{X'',x_1} \wedge_{\mathbb{C}^{*}} T^{*}_{X'',x_2} .
\]

On trouve ainsi un foncteur de

a) catégorie des courbes hol.\ compactes à bord $X$ \uncertain{(quasi-vides)},
de type $g,\nu$, avec disques holomorphes (mais comme flèches les iso)
\note{« $X$ » est ajouté au-dessus de « à bord », et « (quasi-vides) » en
interligne au-dessus de « catégorie des » ; la lecture de ce dernier mot est
très incertaine.}

vers

b) Catégorie des courbes \uncertain{$\mathbb{C}$-alg.}\ stables de type
$(g,\nu)$ $X'''$, munies 1°) d'\struck{une} « structure de
\uncertain{propagation} normale » en ses pts singuliers, et 2°) d'un
\struck{structure de multidisques} \uncertain{syst.}\ admissible de disques
autour des pts marqués.

\page{75}
\note{pagination de l'auteur : 12.}

Au niveau des multiplicités modulaires \marginal{(analytiques réelles)} on
trouve sûrement un morphisme, pour celles-ci
\[
  \mathrm{MDB}_{g,\nu} \longrightarrow \mathrm{M_{st}ND}_{g,\nu}
\]
\note{au-dessus de la flèche, un premier $\mathrm{MDB}_{g,\nu}$ est biffé ;
le second membre porte une surcharge biffée entre le $\mathrm{N}$ et le
$\mathrm{D}$.}
où $\mathrm{MDB}_{g,\nu}$ désigne la multiplicité modulaire (M) pour les
surfaces holomorphes compactes « disquées » (D) à bord (B) de type
$g,\nu$, et $\mathrm{M_{st}ND}_{g,\nu}$ celle pour les courbes algébriques
stables (de type loc.\ constant) \add{de type $g,\nu$} à structure de
propagation normale (N) et à système \uncertain{admissible} de disques (D)
autour des pts marqués.
\note{« de type $g,\nu$ » est en interligne, au-dessus de « (de type loc.\
constant) ».}
Cette flèche est compatible avec les flèches canoniques vers la cat.\ des
types combinatoires \ill{} de tels objets (diagrammes MD (Mumford-Deligne)
« modulaires ») (\struck{indiquant} \uncertain{marquant} une pondération
des sommets) à structures toriques \ill{} aux arêtes… \uncertain{données}
\ill{} diagrammes connexes)
\note{un mot en interligne au-dessus de « combinatoires », terminé en
« -iques », n'est pas lu.}

\page{76}

\begin{tikzcd}
  \mathrm{MDB}_{g,\nu} \arrow[rr] \arrow[dr] & & \mathrm{M_{st}ND}_{g,\nu} \arrow[dl] \\
  & \mathrm{mdt}_{g,\nu} &
\end{tikzcd}

(dans $\mathrm{mdt}$, la lettre $t$ rappelle qu'on a inclus des structures
toriques dans le type combinatoire). Sûrement aussi, si on
\struck{définit} \uncertain{considère} $\mathrm{MD}_{g,\nu}$ (sans bord, ou
à bord « générique » réduit à un ens.\ de $\nu$ pts) et
$\mathrm{M_{st}N}_{g,\nu}$ (sans structure \uncertain{multi-}disquée) on
aura

\begin{tikzcd}
  \mathrm{MD}_{g,\nu} \arrow[rr] \arrow[dr] & & \mathrm{M_{st}N}_{g,\nu} \arrow[dl] \\
  & \mathrm{md}_{g,\nu} &
\end{tikzcd}

cette fois un morphisme de \struck{de} multiplicités \struck{analytiques}
$\mathbb{C}$-algébriques. Dans le cas 1), les dimensions (réelles) \ill{}
sont celles de $\mathbb{C}^{3g-3+\nu} \times \mathbb{R}^{\nu}$ \ill{} les
dim.\ complexes, \uncertain{dans} le cas 2), celle de
$\mathbb{C}^{3g-3+\nu}$.
\note{la fin du paragraphe est une correction en interligne, soulignée, qui
se lit mal ; l'ordre des mots rendu ici est incertain.}

\page{77}
\note{pagination de l'auteur : 13.}

\emph{Question} Ces flèches sont-elles des iso -- et si oui, comment
expliciter un iso en sens inverse ? \uncertain{Sûrt}

Si la détermination conjecturale des syst.\ admissibles de disques sur une
surface holomorphe \uncertain{compacte} (fonctions de type hyperbolique)
était correcte, on trouverait aisément des \struck{\ill{}} flèches en sens
inverse ; (au niveau plus naïf des foncteurs entre les catégories-tests de
ces multiplicités) on \struck{fournirait} \uncertain{aurait} ceci : à
l'aide d'un \struck{\ill{}}
\[
  u_i \in T^{*}_{x'_i} \wedge_{\mathbb{C}^{*}} T^{*}_{x''_i}
\]
\uncertain{ou ce qui revient au même} un
$w_i \,(= u_i^{-1}) \in T^{\vee *}_{x'_i} \wedge T^{\vee}_{x''_i}$,
\struck{i.e.\ un isomorphisme} \uncertain{i.e.} un accouplement dualisant
\[
  T_{x'_i} \otimes T_{x''_i} \longrightarrow \mathbb{C}
\]
\note{sous cet accouplement, une première version est biffée :
$T_{a'_i} \xrightarrow{\ \sim\ } T_{a''_i}$ ; les $x'_i, x''_i$ deviennent
$a'_i, a''_i$ dans les lignes qui suivent.}
\struck{choisissant} choisissant un \struck{\ill{}} disque conforme $D'_i$
dans $T_{a'_i}$, \struck{\ill{}} \uncertain{d'où} disque conforme
\uncertain{polaire} $D''_i$ dans $T_{a''_i}$, on trouverait \struck{des}
pour ce syst.\ $(D'_i)$ $(D''_i)$ des réalisations de ces disques comme

\page{78}

plongés dans $X''$ normalisé de $X'''$ \struck{\ill{}} \ill{} part, et on
trouverait des antiiso entre les bords \struck{de ces} de $D'_i$, $D''_i$
-- d'où possibilité de recoller pour trouver une $X$ surface hol.\ à bord,
si on a aussi enlevé des disques ouverts $\mathring{D}^{*}_{\alpha}$
correspondant aux pts marqués $x_\alpha$ (si on a une surface à pts
marqués), disons par des cercles $\Gamma_i$ qui ne sont peut-être pas (??
\struck{\ill{}} à examiner de plus près) \uncertain{des} géodésiques, mais
définissant des cercles géodésiques uniques, par les principes de
Thurston.
\note{« peut-être pas » est en interligne ; « (?? à examiner de plus
près) » est ajouté au-dessus d'un début de parenthèse biffé.}
Il faudrait voir que ce résultat obtenu ne dépend pas du choix des paires
$(D'_i, D''_i)$ mutuell\supplied{emen}t \uncertain{polaires}, il y a un argument
heuristique pour \uncertain{cela} tout à fait sérieux, consistant à dire
qu'une petite variation de cette

\page{79}
\note{pagination de l'auteur : 14.}

paire ne change pas $X$ ; car si on remplace $D'_i$ par $\lambda D'_i$,
donc $D''_i$ par $\frac{1}{\lambda} D''_i$, avec
$\lambda \in \mathbb{R}^{*+}$ voisin de 1, les \struck{trois} relations
vaudront pour les disques obtenus dans $X''$ \uncertain{normalisé} (c'est
probablement tout à fait faux !), et on obtiendrait la même surface $X$,
mais les $\Gamma_i$ remplacés par des cercles « voisins » $\Gamma'_i$.

Le cas-clef à comprendre sans doute est celui où $(g,\nu) = (0,4)$, i.e.\
le \uncertain{recollement} de deux « pantalons » de Thurston,
\uncertain{de} deux \struck{\ill{}} \uncertain{données} $\Sigma'$,
$\Sigma''$ à \uncertain{trois} éléments, fixés, la fonction de

\begin{quote}
$\{$pantalons de type combinatoire $\Sigma' \sqcup \{s\}$,
$\Sigma'' \sqcup \{t\}$, \struck{donnés} recollés suivant $\{s, t\}$, les
données $(l_i)_{i \in \Sigma' \sqcup \Sigma'' \sqcup u} \in
(\mathbb{R}^{*+})^{\Sigma' \sqcup \Sigma'' \sqcup u}$, et
$\theta \in \mathbb{U}_{\Sigma' \wedge \Sigma''}$ qui fixe le
recollement$\}$
\end{quote}

vers

\page{80}

\begin{quote}
$\{$paires de pantalons \struck{\ill{}} \uncertain{généralisés} $X'$,
$X''$ \uncertain{à bord} indexés par $\Sigma' \sqcup \{s\}$,
$\Sigma'' \sqcup \{t\}$, avec $s, t$ correspondant aux \struck{bords}
composantes ponctuelles du bord, avec les périmètres
$(\lambda_i)_{i \in \Sigma' \sqcup \Sigma''} \in
\mathbb{R}^{\Sigma' \sqcup \Sigma''}$, \struck{et} et un isom
d'identification des $\mathbb{C}^{*}_{\pm 1}(\Sigma' \wedge \Sigma'')\}$
\end{quote}
\note{chacun des deux ensembles est marqué d'une grande accolade en marge
gauche ; l'indice $u$ et le groupe $\mathbb{U}_{\Sigma' \wedge \Sigma''}$
sont lus sous réserve.}

Une fois ce cas compris, et revenant alors au cas général, il faudrait
déjà, pour un type de diagramme donné -- \uncertain{incarné} par un
diagramme MD dans $\mathrm{md}_{g,\nu}$, comprendre l'application sur les
groupes fondamentaux induite par
\[
  \mathrm{MDB} \longrightarrow \mathrm{M_{st}N}_{g,\nu} ,
\]
\marginal{(*)} celui du deuxième membre étant \ill{} complété\supplied{s}, un
\ill{} de groupes de Teichmüller spéciaux -- \emph{mais il doit être pareil
pour le premier} ! Également, il ne devrait pas être bien dur de vérifier
si
\note{la page s'arrête ici, au milieu de la phrase.}

\end{document}
